\subsection*{الف}
با توجه به
$h_i=\tanh(w h_{i-1})$
داریم
$-1<h_i<1$
پس
$0<1-h_i^2\le 1$.
بنابراین
\begin{flalign*}
    \gamma_i = w(1-h_i^2)
    \;\;\Rightarrow\;\;
    |\gamma_i| = |w|\,(1-h_i^2),
\end{flalign*}
و در نتیجه
\begin{flalign*}
    0 \le |\gamma_i| \le |w|. &&
\end{flalign*}

\subsection*{ب}
از رابطه‌ی گرادیان، ترم اصلی به صورت
\begin{flalign*}
    w^{T-t}\prod_{j=t+1}^{T}(1-h_j^2)
    = \prod_{j=t+1}^{T} |w|\,(1-h_j^2)
    = \prod_{j=t+1}^{T} |\gamma_j|
\end{flalign*}
ظاهر می‌شود. بنابراین شرط محو شدن گرادیان این است که
\begin{flalign*}
    \prod_{j=t+1}^{T} |\gamma_j| \longrightarrow 0 \quad (T\to\infty).
\end{flalign*}
یک شرط کافی:
\begin{flalign*}
    \exists\,\rho<1:\; |\gamma_j|\le \rho \;\;\text{برای همه‌ی}\; j\ge j_0. &&
\end{flalign*}
به طور خاص اگر
$|w|<1$
باشد، همواره گرادیان محو می‌شود.

\subsection*{ج}
انفجار گرادیان زمانی رخ می‌دهد که
\begin{flalign*}
    \prod_{j=t+1}^{T} |\gamma_j| \longrightarrow \infty \quad (T\to\infty).
\end{flalign*}
یک شرط کافی:
\begin{flalign*}
    \exists\,\rho>1:\; |\gamma_j|\ge \rho \;\;\text{برای همه‌ی}\; j\ge j_0. &&
\end{flalign*}
از آنجا که
$|\gamma_j|=|w|(1-h_j^2)\le |w|$،
برای انفجار لازم است
$|w|>1$
و حالت‌های پنهان اشباع نشده باشند، یعنی
$(1-h_j^2)$
به اندازه‌ی کافی بزرگ بماند.